% META: {"id": "TCOMPL-BAYES-EV-A", "chapitre": "TCOMPL-INFERENCE-BAYESIENNE", "type_objet": "evaluation", "version": "A", "duree_min": 45, "bareme_total": 20, "capacites_codes": ["C1", "C2", "C3"], "competences": ["calculer", "raisonner"], "mode_creation": "ex_nihilo", "sources_inspiration": [], "status": "generated"}
% BEGIN-VERIFY
% from sympy import Rational
% PA = Rational(3, 100)
% PA_B = Rational(85, 100)
% PnotA_B = Rational(4, 100)
% PB = PA*PA_B + (1-PA)*PnotA_B
% assert PB == Rational(643, 10000)
% PB_A = PA_B*PA/PB
% assert PB_A == Rational(255, 643)
% assert abs(float(PB_A) - 0.3966) < 0.001
% END-VERIFY

%---------------------------------------------------------------
% Duree : 45 minutes — Bareme : 20 points
%---------------------------------------------------------------

\begin{evaluation}{TCOMPL-BAYES-EV-A}{45}

\begin{exercice}{TCOMPL-BAYES-EV-A-EX1}{2}{40}
\textbf{Exercice unique} (20 points) --- \textit{Capacités C1, C2, C3}

\medskip

Dans un pays, $3\%$ de la population est porteuse d'un anticorps particulier ($P(A)=0{,}03$). Un test le détecte avec une probabilité de $85\%$ quand il est présent, et donne un résultat positif à tort avec une probabilité de $4\%$ chez les personnes non porteuses.
\begin{enumerate}
  \item (4 pts) Rappeler la formule de Bayes reliant $P_B(A)$, $P_A(B)$, $P(A)$ et $P(B)$.
  \item (6 pts) Calculer $P(B)$, la probabilité qu'un individu pris au hasard ait un test positif.
  \item (6 pts) En déduire $P_B(A)$.
  \item (4 pts) Expliquer pourquoi $P_B(A)$ et $P_A(B)=0{,}85$ sont si différents, en une phrase.
\end{enumerate}
\end{exercice}

\end{evaluation}
